% Generated by tools/build_new_dists_anthology_sections.py; do not edit by hand.

\section{Finite local-ring permutation actions}\label{proof:new-dist:local_ring_actions}

\resultbox{\textbf{Canonical testing artifacts.}\par\smallskip Exact backend: \path{lambda_distributions/dists/local_ring_permutation_actions/verification.py}.\par Regression: \path{lambda_distributions/dists/local_ring_permutation_actions/test_verification.py}.\par Marimo: \path{marimo_notebooks/local_ring_actions.py}.}

\subsection*{Scope and outcome}

Let $R_a=\calO/\pi^a\calO$ be a finite chain ring with residue field
$\F_q$.  A subgroup of $M_n(R_a)$ is not canonically a complex matrix group;
the representations studied here are instead the honest complex permutation
modules $\C[X]$ on primitive vectors, free projective points, free summands,
flags, formed-space submodules, and congruence adjoint quotients.  We prove the
fixed-point, cycle, and $\sigma$-MGF formulas, the low-degree orbit formulas,
and fixed-level rank stability.  An executable marimo notebook independently
constructs representative groups over $\mathbb Z/4$ and $\mathbb Z/9$ and
compares direct configuration orbits with Smith-kernel, homogeneous-space,
M\"obius-cycle, and adjoint-kernel calculations.  All reported checks pass
exactly.

\subsection{Representation convention and universal permutation theorem}

Fix a finite group $H$ acting on a finite set $X$, and put $V=\C[X]$.  For
$h\in H$, write
\[
 F_d(h;X)=|\Fix_X(h^d)|,
 \qquad c_r(h;X)=\#\{\text{$r$-cycles of $h$ on $X$}\}.
\]
A point on an $r$-cycle is fixed by $h^d$ exactly when $r\mid d$, so
\[
 F_d(h;X)=\sum_{r\mid d}r c_r(h;X).
\]
M\"obius inversion proves
\begin{equation}\label{nd:local_ring_actions:eq:mobius}
 \boxed{c_r(h;X)=\frac1r\sum_{d\mid r}\mu(r/d)F_d(h;X).}
\end{equation}

The permutation matrix of an $r$-cycle has all $r$th roots of unity as its
eigenvalues.  Its characteristic factor is therefore $1-t^r$.  Reynolds
averaging and the symmetric-power identity give the exact series
\begin{equation}\label{nd:local_ring_actions:eq:universal}
 \boxed{\MGF_{H,X}(t_1,t_2,\ldots)=\frac1{|H|}\sum_{h\in H}
 \prod_{i\ge1}\prod_{r\ge1}(1-t_i^r)^{-c_r(h;X)}.}
\end{equation}
Indeed, for a partition or composition $\tau=(\tau_1,\ldots,\tau_s)$,
\[
 [m_\tau]\MGF_{H,X}
 =\dim\left(\bigotimes_{j=1}^s\Sym^{\tau_j}\C[X]\right)^H.
\]
A monomial basis for $\Sym^d\C[X]$ is indexed by $d$-element multisets in
$X$.  Orbit sums are a basis of the invariants, proving the independent
interpretation
\begin{equation}\label{nd:local_ring_actions:eq:orbits}
 \boxed{[m_\tau]\MGF_{H,X}
 =\#\left(H\backslash\prod_j\Sym^{\tau_j}X\right).}
\end{equation}
Equations \eqref{nd:local_ring_actions:eq:mobius}--\eqref{nd:local_ring_actions:eq:orbits} reduce every family in this
note to a fixed-point problem on a finite set.

\subsection[GLn over Ra on primitive vectors]{$\GL_n(R_a)$ on primitive vectors}

Let
\[
 R_a=\calO/\pi^a\calO,\qquad R_a/\pi R_a\simeq\F_q,
 \qquad X^{\rm prim}_{n,a}=R_a^n\setminus\pi R_a^n.
\]
For $A\in M_n(R_a)$, let $0\le\nu_1(A),\ldots,\nu_n(A)\le a$ be its
truncated Smith exponents; a zero Smith entry receives exponent $a$.  Define
\begin{equation}\label{nd:local_ring_actions:eq:Kb}
 K_b(A)=q^{\sum_{j=1}^n\min(\nu_j(A),b)}.
\end{equation}
Smith reduction shows that $K_b(A)$ is the size of the kernel of $A$ after
reduction to $R_b$.

\subsubsection{Fixed points and exact series}
A primitive vector fixed by $h^d$ lies in $\ker(h^d-I)$ but not in
$\pi R_a^n$.  Multiplication by $\pi$ identifies $R_{a-1}^n$ with
$\pi R_a^n$, and under this identification the second kernel has size
$K_{a-1}(h^d-I)$.  Hence
\begin{equation}\label{nd:local_ring_actions:eq:primitive-fixed}
 \boxed{F_d^{\rm prim}(h)=K_a(h^d-I)-K_{a-1}(h^d-I).}
\end{equation}
Substitution in \eqref{nd:local_ring_actions:eq:mobius} and \eqref{nd:local_ring_actions:eq:universal} proves
\begin{equation}\label{nd:local_ring_actions:eq:primitive-smg}
 \boxed{\MGF^{\rm prim}_{n,a}=\frac1{|\GL_n(R_a)|}\sum_h
 \prod_{i,r\ge1}(1-t_i^r)^{-\frac1r\sum_{d\mid r}\mu(r/d)
 [K_a(h^d-I)-K_{a-1}(h^d-I)]}.}
\end{equation}
The identity element in \eqref{nd:local_ring_actions:eq:primitive-fixed} also gives
$|X^{\rm prim}_{n,a}|=q^{an}-q^{(a-1)n}$.

\subsubsection{Ordered-pair coefficient}
For $n\ge2$, transitivity allows the first primitive vector to be fixed as
$e_1$.  Write the second as $w=xe_1+y$.  Its component $y$ in a chosen
complement has valuation $r\in\{0,\ldots,a\}$, with $r=a$ meaning $y=0$.
The stabilizer of $e_1$ moves $y$ to $\pi^r e_2$ and changes $x$ by an
element of $\pi^rR_a$.  For $r=0$ there is one orbit.  For $1\le r\le a$,
primitivity forces the residue of $x$ modulo $\pi^r$ to be a unit, giving
$q^r-q^{r-1}$ orbits.  Consequently
\begin{equation}\label{nd:local_ring_actions:eq:primitive-pair}
 \boxed{[m_{(1,1)}]\MGF^{\rm prim}_{n,a}
 =1+\sum_{r=1}^a(q^r-q^{r-1})=q^a.}
\end{equation}

\subsection[GLn over Ra on projective points]{$\GL_n(R_a)$ on projective points}

Let $X^{\rm proj}_{n,a}=\mathbf P^{n-1}(R_a)$ be the free rank-one direct
summands.  If $L=R_av$ is fixed by $h^d$, then $h^dv=uv$ for a unique
$u\in R_a^\times$.  Conversely each primitive $u$-eigenvector generates a
fixed line, and every line has $|R_a^\times|$ primitive generators.  Thus
\begin{equation}\label{nd:local_ring_actions:eq:projective-fixed}
 \boxed{F_d^{\rm proj}(h)=\frac1{|R_a^\times|}\sum_{u\in R_a^\times}
 \left[K_a(h^d-uI)-K_{a-1}(h^d-uI)\right].}
\end{equation}
Together with \eqref{nd:local_ring_actions:eq:mobius} this yields
\begin{equation}\label{nd:local_ring_actions:eq:projective-smg}
 \boxed{\MGF^{\rm proj}_{n,a}=\frac1{|\GL_n(R_a)|}\sum_h
 \prod_{i,r\ge1}(1-t_i^r)^{-\frac1r\sum_{d\mid r}\mu(r/d)
 F_d^{\rm proj}(h)}.}
\end{equation}
Counting primitive generators gives
\[
 |\mathbf P^{n-1}(R_a)|
 =q^{(a-1)(n-1)}\frac{q^n-1}{q-1}.
\]

For $n\ge2$, after fixing the first line, every second line has a unique
relative-position parameter
\[
 R_a(e_1+\pi^r e_2),\qquad 0\le r\le a,
\]
up to the stabilizer.  The parameter is unchanged when the two lines are
interchanged.  It classifies both ordered pairs and two-element multisets,
including equality at $r=a$.  Therefore
\begin{equation}\label{nd:local_ring_actions:eq:projective-pair}
 \boxed{[m_{(1,1)}]\MGF^{\rm proj}_{n,a}
 =[m_{(2)}]\MGF^{\rm proj}_{n,a}=a+1.}
\end{equation}
For $n=2$ this is also the complete-flag action.  At $a=2$ its three
orbitals already exceed the field-level Weyl count $|S_2|=2$, concretely
exhibiting the extra valuation parameter in local-ring Bruhat theory.

\subsection{Free summands, flags, and homogeneous spaces}

Let $H=\GL_n(R_a)$ and let $P_{k,n-k}$ stabilize the standard free
rank-$k$ summand.  For any finite homogeneous space $H/P$,
\begin{equation}\label{nd:local_ring_actions:eq:coset-fixed}
 |\Fix_{H/P}(x)|=\frac{|C_H(x)|\,|x^H\cap P|}{|P|}.
\end{equation}
To prove this, count $y\in H$ satisfying $y^{-1}xy\in P$.  Every fixed coset
has exactly $|P|$ such representatives.  On the other hand, each element of
$x^H\cap P$ has exactly $|C_H(x)|$ conjugators.

For $X^{(k)}_{n,a}=\Gr_k^{\rm free}(R_a^n)=H/P_{k,n-k}$, equations
\eqref{nd:local_ring_actions:eq:coset-fixed}, \eqref{nd:local_ring_actions:eq:mobius}, and \eqref{nd:local_ring_actions:eq:universal} give
\begin{align}
 F_d^{(k)}(h)
 &=\frac{|C_H(h^d)|\,|(h^d)^H\cap P_{k,n-k}|}{|P_{k,n-k}|},\label{nd:local_ring_actions:eq:grass-fixed}\\
 \MGF^{(k)}_{n,a}
 &=\frac1{|H|}\sum_h\prod_{i,r\ge1}
 (1-t_i^r)^{-\frac1r\sum_{d\mid r}\mu(r/d)F_d^{(k)}(h)}.
 \label{nd:local_ring_actions:eq:grass-smg}
\end{align}
Reduction modulo $\pi$ and graph lifting over a fixed complement yield
\[
 |X^{(k)}_{n,a}|=q^{(a-1)k(n-k)}\qbinom{n}{k}.
\]

\subsubsection{Pair coefficient for free summands}
Assume $n\ge2k$ and fix one summand $U$.  For a second summand $W$, apply
Smith reduction to the projection $W\to R_a^n/U$.  Its uniquely determined
exponents satisfy
$0\le\lambda_1\le\cdots\le\lambda_k\le a$.  Changes of bases in $W$, in
$U$, and in a complementary $k$-summand put the pair in the form
\[
 U=\langle e_1,\ldots,e_k\rangle,\qquad
 W=\langle e_i+\pi^{\lambda_i}e_{k+i}:1\le i\le k\rangle.
\]
Parabolic shears remove the remaining components, so the Smith exponents are
also a complete invariant.  There are as many nondecreasing $k$-tuples from
$\{0,\ldots,a\}$ as multisets of size $k$ from an $(a+1)$-element set:
\begin{equation}\label{nd:local_ring_actions:eq:grass-pair}
 \boxed{[m_{(1,1)}]\MGF^{(k)}_{n,a}
 =[m_{(2)}]\MGF^{(k)}_{n,a}=\binom{a+k}{k}.}
\end{equation}
At $a=1$ this reduces to the $k+1$ field-level intersection types.

For a partial flag shape $\mathbf d=(d_1<\cdots<d_s)$, put
$X_{\mathbf d,n,a}=H/P_{\mathbf d,n}$.  The same proof gives
\begin{equation}\label{nd:local_ring_actions:eq:flag}
 F_d(h;X_{\mathbf d,n,a})
 =\frac{|C_H(h^d)|\,|(h^d)^H\cap P_{\mathbf d,n}|}{|P_{\mathbf d,n}|},
\end{equation}
followed by \eqref{nd:local_ring_actions:eq:mobius} and \eqref{nd:local_ring_actions:eq:universal}.  Unlike the field
case, $P_{\mathbf d,n}\backslash H/P_{\mathbf d,n}$ generally retains
valuation data and is not indexed only by the Weyl group.

\subsection{Fixed-level stability as rank grows}

Fix $a$, take $Q\le\GL_d(R_a)$, and define
\[
 X_n(Q)=\operatorname{SplitInj}(R_a^d,R_a^n)/Q.
\]
Primitive vectors, projective points, rank-$k$ summands, and fixed-shape flags
arise from $d=1,Q=1$; $d=1,Q=R_a^\times$; $d=k,Q=\GL_k(R_a)$; and
$d=d_s,Q=P_{\mathbf d}(R_a)$, respectively.

Let $D=|\tau|$.  Choose a split frame for each of the $D$ objects and
concatenate the resulting blocks into an $n\times dD$ matrix $A$.  Let $L(A)$
be its row submodule in $R_a^{dD}$.  Every block projection $L(A)\to R_a^d$
is surjective.  Write
\[
 \mathcal L_{\tau,d}(R_a)=\{L\le R_a^{dD}:\text{every block projection is
 surjective}\},\qquad
 \Gamma_{\tau,Q}=Q^D\rtimes\prod_jS_{\tau_j}.
\]

Two matrices with the same row submodule are two generating $n$-tuples of
the same finite $R_a$-module.  Over a local ring, $\GL_n(R_a)$ is transitive
on such generating tuples once $n$ is at least the minimal number of
generators: reduce modulo $\pi$, move a basis into the first coordinates,
eliminate the other generators, and lift the remaining change of basis
through the invertible group $I+\pi M_n(R_a)$.  Every submodule of
$R_a^{dD}$ needs at most $dD$ generators.  The block changes and within-color
reorderings are precisely $\Gamma_{\tau,Q}$.  Hence
\begin{equation}\label{nd:local_ring_actions:eq:stable}
 \boxed{[m_\tau]\MGF_{\GL_n(R_a),X_n(Q)}
 =\#\bigl(\Gamma_{\tau,Q}\backslash\mathcal L_{\tau,d}(R_a)\bigr),
 \qquad n\ge dD.}
\end{equation}
The right side is independent of $n$.  For $\SL_n(R_a)$, the same coefficient
agrees in the safe range $n>dD$: an unused coordinate corrects the determinant
without moving the configuration.  The conclusion passes to the corresponding
projective quotients after factoring the action kernel.

\clearpage
\subsection{Symplectic and orthogonal towers}
\thispagestyle{fancy}

Let $G_n$ be a finite symplectic or orthogonal group over $R_a$ and let $X_n$
be a homogeneous orbit of split totally isotropic or singular free submodules
or flags of fixed shape.  The exact statement needs no new averaging:
\begin{equation}\label{nd:local_ring_actions:eq:formed}
 \boxed{\MGF_{G_n,X_n}=\frac1{|G_n|}\sum_{g\in G_n}\prod_{i,r\ge1}
 (1-t_i^r)^{-\frac1r\sum_{d\mid r}\mu(r/d)|\Fix_{X_n}(g^d)|}.}
\end{equation}
When $X_n=G_n/P_n$, equation \eqref{nd:local_ring_actions:eq:coset-fixed} computes its fixed
points.  For fixed $a$, shape, and coefficient degree, each configuration is
supported on a formed submodule of bounded rank.  Under the standard split,
unimodular hypotheses, Witt extension identifies two ambient configurations
exactly when their bounded supports are isometric.  Only finitely many support
types exist over fixed $R_a$, so the coefficient is eventually independent of
Witt rank.  In residue characteristic two, the quadratic form itself must be
retained; its polar bilinear form alone is not a complete orthogonal invariant.

\subsection{Reductive group schemes and adjoint congruence quotients}

Let $\mathcal G$ be a reductive group scheme, $G_a=\mathcal G(R_a)$, and let
$\rho_a$ be any honest complex representation with character $\chi_a$.
Expanding each inverse determinant logarithm gives the class-compressed master
formula
\begin{equation}\label{nd:local_ring_actions:eq:character}
 \boxed{\MGF_{G_a,\rho_a}=\sum_{[g]\subseteq G_a}\frac1{|C_{G_a}(g)|}
 \exp\!\left(\sum_{r\ge1}\frac{\chi_a(g^r)}r p_r\right).}
\end{equation}
If $X_a=G_a/P_a$ is a single parabolic orbit, combine
\eqref{nd:local_ring_actions:eq:coset-fixed} with \eqref{nd:local_ring_actions:eq:universal}.  If
$(\mathcal G/\mathcal P)(R_a)$ has several $G_a$-orbits, the formula must be
applied orbit by orbit; silently treating the set as one coset space would be
incorrect.

Now let $c\le a$ and
$\mathfrak g_c=\Lie(\mathcal G)(R_c)$.  Let $G_a$ act on this finite set by
reduction modulo $\pi^c$ and the adjoint action.  Then
\begin{equation}\label{nd:local_ring_actions:eq:adjoint-fixed}
 \boxed{F_d^{\Ad,c}(g)=\left|\ker\bigl(\Ad(\bar g^d)-I:
 \mathfrak g_c\to\mathfrak g_c\bigr)\right|.}
\end{equation}
If the Smith exponents of the displayed operator are
$\eta_1,\ldots,\eta_N\in\{0,\ldots,c\}$, the kernel has size
$q^{\eta_1+\cdots+\eta_N}$.  Therefore
\begin{equation}\label{nd:local_ring_actions:eq:adjoint-smg}
 \boxed{\MGF_{G_a,\mathfrak g_c}=\frac1{|G_a|}\sum_g\prod_{i,r\ge1}
 (1-t_i^r)^{-\frac1r\sum_{d\mid r}\mu(r/d)
 q^{\sum_j\eta_j(g^d)}}.}
\end{equation}

\subsection{Level growth and the correct asymptotic object}

The rank-stable degree-two coefficients still grow with the level:
\[
 \begin{array}{c|c}
 \text{action}&[m_{(1,1)}]\text{ in stable rank}\\ \hline
 \text{primitive vectors}&q^a\\
 \text{projective points}&a+1\\
 \text{rank-$k$ free summands}&\binom{a+k}{k}.
 \end{array}
\]
Thus no unnormalized coefficientwise limit exists as $a\to\infty$ for these
actions.  The natural replacement is the level-generating series.  With the
formal level-zero convention,
\begin{equation}\label{nd:local_ring_actions:eq:level-series}
 \sum_{a\ge0}q^a u^a=\frac1{1-qu},\qquad
 \sum_{a\ge0}(a+1)u^a=\frac1{(1-u)^2},\qquad
 \sum_{a\ge0}\binom{a+k}{k}u^a=\frac1{(1-u)^{k+1}}.
\end{equation}
For actual rings indexed from $a=1$, subtract the constant term.  The growth
is exponential for primitive pairs and polynomial of degree $1$ or $k$ for
projective points or rank-$k$ summands.

\subsection{Independent computational verification}

The accompanying executable files are
\begin{center}
\path{verification.py},\qquad
\path{marimo_notebooks/local_ring_actions.py},\qquad
\path{test_verification.py}.
\end{center}
The marimo notebook is separated into four group/action sections.  It performs
the following independent checks.
\begin{enumerate}
\item Enumerate matrices over $\mathbb Z/p^a$ and retain exactly those whose
determinant is a unit; retain determinant one for $\Sp_2=\SL_2$.
\item Construct primitive vectors and free rank-one summands as literal finite
sets, then induce the matrix permutations.
\item Count orbits directly on tuples of multisets, without using characters,
fixed-point formulas, or cycle indices.
\item Separately compute the cycle-index coefficient in
\eqref{nd:local_ring_actions:eq:universal}; reconstruct cycles from fixed points of powers using
\eqref{nd:local_ring_actions:eq:mobius}.
\item Recover truncated Smith exponents from exact kernel sizes over every
$R_b$ and compare \eqref{nd:local_ring_actions:eq:primitive-fixed} and
\eqref{nd:local_ring_actions:eq:projective-fixed} with literal fixed points.
\item Check \eqref{nd:local_ring_actions:eq:coset-fixed} from centralizers and conjugacy-class
intersections with the projective parabolic.
\item Build all rank-two free summands of $(\mathbb Z/4)^4$ as unique graph
lifts of planes in $\F_2^4$, then compute stabilizer orbits from parabolic
generators.
\item Construct the adjoint action by explicit conjugation and compare every
fixed-point count with \eqref{nd:local_ring_actions:eq:adjoint-fixed}.
\end{enumerate}

\subsubsection{Exact result ledger}

\begin{center}
\small
\begin{tabular}{@{}lrrrr@{}}
\toprule
Action & $|G|$ & $|X|$ & $\tau$ & direct/formula\\
\midrule
$\GL_2(\mathbb Z/4)$, primitive & 96&12&$(1,1)$&$4/4$\\
$\GL_2(\mathbb Z/4)$, projective &96&6&$(2)$&$3/3$\\
&&&$(1,1)$&$3/3$\\
$\GL_2(\mathbb Z/9)$, primitive &3888&72&$(1,1)$&$9/9$\\
$\GL_2(\mathbb Z/9)$, projective &3888&12&$(2)$&$3/3$\\
&&&$(1,1)$&$3/3$\\
$\Sp_2(\mathbb Z/4)$, isotropic lines &48&6&$(1,1)$&$3/3$\\
$\Sp_4(2)$, isotropic points &720&15&$(2),(1,1)$&$3/3$ each\\
$O_4^+(2)$, singular points &72&9&$(2),(1,1)$&$3/3$ each\\
$\GL_2(\mathbb Z/4)$, $M_2(\F_2)$ adjoint &96&16&$(1),(2),(1,1)$
&$6/6,34/34,56/56$\\
\bottomrule
\end{tabular}
\end{center}

The $\GL_2(\mathbb Z/4)$ Smith, homogeneous-space, and M\"obius checks cover
all $96$ matrices and every relevant power.  For $\GL_2(\mathbb Z/9)$ the
orbit and cycle averages use all $3888$ matrices, while the more expensive
pointwise comparison uses $32$ deterministic representatives.  The rank-two
summand construction contains exactly
\[
 2^{(2-1)2(4-2)}\left.\qbinom{4}{2}\right|_{q=2}=16\cdot35=560
\]
summands.  The stabilizer has exactly six orbits, with Smith types
\[
 (0,0),(0,1),(0,2),(1,1),(1,2),(2,2),
\]
which is the complete list and agrees with $\binom{2+2}{2}=6$.

\resultbox{\textbf{Verification outcome.}  Ten newly tabulated coefficient
comparisons pass exactly.  All fixed-point/cycle reconstructions pass in the
stated exhaustive or representative ranges.  Four automated tests complete
successfully, and the marimo notebook passes its static notebook check.}

\subsection{Scope and reproducibility}

The proof applies to finite chain rings with residue cardinality $q$.  The
executable constructors specialize to $R_a=\mathbb Z/p^a\mathbb Z$, so their
residue cardinality is $q=p$; this is sufficient to test non-field level
behavior at two distinct primes.  The computations prove no untested
all-parameter statement: the universal Reynolds, Smith, coset-counting, and
bounded-support arguments above provide those proofs.

From the workspace root, run
\begin{verbatim}
marimo run marimo_notebooks/local_ring_actions.py
pytest -q lambda_distributions/dists/local_ring_permutation_actions/test_verification.py
\end{verbatim}
The standalone PDF generated from this source is designed to be included
later, without modification, by the proof anthology's \verb|\proofnote|
mechanism.

\begin{thebibliography}{9}
\bibitem{nd:local_ring_actions:grassmann-residue}
Y.~Guo, J.~Guo, and L.~Huang,
\emph{Vector spaces and Grassmann graphs over residue class rings},
\href{https://arxiv.org/abs/1705.04610}{arXiv:1705.04610}.

\bibitem{nd:local_ring_actions:bruhat-local}
U.~Onn, A.~Prasad, and L.~Vaserstein,
\emph{A note on Bruhat decomposition of $GL(n)$ over local principal ideal
rings},
\href{https://arxiv.org/abs/math/0506094}{arXiv:math/0506094}.

\bibitem{nd:local_ring_actions:witt-extension}
J.-P.~Tignol,
\emph{Witt's extension theorem for quadratic spaces over semiperfect rings},
\href{https://arxiv.org/abs/1408.0522}{arXiv:1408.0522}.
\end{thebibliography}
